\section{Ньютон}

\subsection{Понятие вектора}

\textbf{Вектор} — математический объект, характеризующийся численным значением \textit{(модулем, или длиной)} и направлением в пространстве. Вектор изображается направленным отрезком, где длина отрезка соответствует модулю вектора, а стрелка указывает его направление.

\textbf{Длина (модуль) вектора} $\vec{a} = (a_x, a_y, a_z)$ в трёхмерном декартовом пространстве вычисляется по формуле:
\[
|\vec{a}| = \sqrt{a_x^2 + a_y^2 + a_z^2}.
\]

Свойства длины вектора:
\begin{itemize}
    \item $|\vec{a}| \geq 0$, причём $|\vec{a}| = 0$ тогда и только тогда, когда $\vec{a} = \vec{0}$ (нулевой вектор).
    \item Для любого скаляра $\lambda$: $|\lambda \vec{a}| = |\lambda| \cdot |\vec{a}|$.
    \item Неравенство треугольника: $|\vec{a} + \vec{b}| \leq |\vec{a}| + |\vec{b}|$.
\end{itemize}

\subsection{Операции над векторами}

\subsubsection{Сложение и вычитание}

\textbf{Сложение векторов} — это операция нахождения вектора суммы $\vec{c} = \vec{a} + \vec{b}$:

\begin{enumerate}
    \item \textbf{Геометрический способ:}
    \begin{itemize}
        \item \textit{Правило треугольника:} Вектор $\vec{b}$ совмещается своим началом с концом вектора $\vec{a}$. Вектор суммы $\vec{c}$ соединяет начало $\vec{a}$ с концом $\vec{b}$.
        \item \textit{Правило параллелограмма:} Векторы $\vec{a}$ и $\vec{b}$ совмещаются своими началами. Вектор суммы $\vec{c}$ является диагональю параллелограмма, построенного на этих векторах, выходящей из их общего начала.
    \end{itemize}
    
    \item \textbf{Аналитический способ:}
    Сумма векторов находится как сумма их соответствующих координат в декартовой системе:
    \[
        \vec{c} = (a_x + b_x)\vec{i} + (a_y + b_y)\vec{j} + (a_z + b_z)\vec{k}
    \]
\end{enumerate}

\begin{tblr}{
    colspec = {X[4.5,c]X[8.5,l]},
    vlines,
    hlines,
    row{1} = {font=\bfseries, mode=text},
    row{odd} = {bg=gray!10},
    cells = {m},
    rowsep = 6pt
}
Свойство / Операция & Описание \\
$\displaystyle \vec{a} + \vec{b} = \vec{b} + \vec{a}$ & Коммутативность: результат не зависит от порядка слагаемых \\
$\displaystyle (\vec{a} + \vec{b}) + \vec{d} = \vec{a} + (\vec{b} + \vec{d})$ & Ассоциативность: при сложении трёх и более векторов их можно группировать любым способом \\
$\displaystyle \vec{d} = \vec{a} - \vec{b} = \vec{a} + (-\vec{b})$ & Вычитание: вектор разности $\vec{d}$ направлен от конца вычитаемого ($\vec{b}$) к концу уменьшаемого ($\vec{a}$) \\
\end{tblr}

\subsubsection{Скалярное произведение}
\textbf{Скалярным произведением} двух векторов $\vec{a}$ и $\vec{b}$ называется скалярная величина, равная произведению модулей этих векторов на косинус угла между ними:

\begin{center}
\begin{tblr}{
    colspec = {X[4.5,c]X[8.5,l]},
    vlines,
    hlines,
    row{1} = {font=\bfseries, mode=text},
    row{odd} = {bg=gray!10},
    cells = {m},
    rowsep = 6pt
}
Формула & Описание \\
$\displaystyle \vec{a} \cdot \vec{b} = |\vec{a}| |\vec{b}| \cos \alpha$ & Геометрическое определение; $\alpha$ — угол между векторами \\
$\displaystyle \vec{a} \cdot \vec{b} = a_x b_x + a_y b_y + a_z b_z$ & Алгебраическое определение в декартовых координатах \\
$\displaystyle \vec{a} \cdot \vec{b} = \vec{b} \cdot \vec{a}$ & Коммутативность (перестановочный закон) \\
$\displaystyle \vec{a} \cdot \vec{a} = |\vec{a}|^2$ & Скалярный квадрат вектора равен квадрату его модуля \\
$\displaystyle \vec{a} \cdot \vec{b} = 0 \iff \vec{a} \perp \vec{b}$ & Условие ортогональности: произведение равно нулю, если векторы перпендикулярны \\
\end{tblr}
\end{center}

\subsubsection{Векторное произведение векторов}
\textbf{Векторным произведением} вектора $\vec{a}$ на вектор $\vec{b}$ называется вектор $\vec{c} = [\vec{a}, \vec{b}]$, удовлетворяющий трем условиям:
\begin{enumerate}
    \item Модуль вектора: $|\vec{c}| = |\vec{a}| |\vec{b}| \sin \alpha$.
    \item Вектор $\vec{c}$ перпендикулярен плоскости, в которой лежат векторы $\vec{a}$ и $\vec{b}$.
    \item Векторы $\vec{a}, \vec{b}, \vec{c}$ образуют \textit{правую тройку} (направление определяется правилом правой руки/буравчика).
\end{enumerate}

\begin{center}
\begin{tblr}{
    colspec = {X[4.5,c]X[8.5,l]},
    vlines,
    hlines,
    row{1} = {font=\bfseries, mode=text},
    row{odd} = {bg=gray!10},
    cells = {m},
    rowsep = 6pt
}
Свойство & Описание \\
$\displaystyle [\vec{a}, \vec{b}] = -[\vec{b}, \vec{a}]$ & Антикоммутативность: при перестановке множителей вектор меняет направление на противоположное \\
$\displaystyle [\vec{a}, \vec{b}] = 0 \iff \vec{a} \parallel \vec{b}$ & Условие коллинеарности: произведение равно нулю, если векторы параллельны \\
$\displaystyle S_{\text{пар}} = |[\vec{a}, \vec{b}]|$ & Геометрический смысл: модуль векторного произведения равен площади параллелограмма, построенного на векторах \\
\end{tblr}
\end{center}

\subsection{Кинематика материальной точки}

\begin{tblr}{
    colspec = {X[4,c]X[8,l]},
    vlines,
    hlines,
    row{1} = {font=\bfseries, mode=text},
    row{odd} = {bg=gray!10},
    cells = {m},
    rowsep = 6pt
}
Величина & Описание и связь \\
$\displaystyle \Delta \vec{r} = \vec{r}_2 - \vec{r}_1$ & \textit{Перемещение:} вектор, соединяющий начальное и конечное положение точки \\
$\displaystyle s = \int |v| dt$ & \textit{Путь:} скалярная величина, равная длине траектории; $s \ge |\Delta \vec{r}|$ \\
$\displaystyle \vec{v} = \frac{d\vec{r}}{dt}$ & \textit{Мгновенная скорость:} производная радиус-вектора по времени, направлена по касательной к траектории \\
$\displaystyle \langle v \rangle = \frac{s}{t}$ & \textit{Средняя скорость:} скалярная величина, равная отношению всего пройденного пути ко всему времени движения \\
$\displaystyle \vec{a} = \frac{d\vec{v}}{dt} = \frac{d^2\vec{r}}{dt^2}$ & \textit{Ускорение:} величина, характеризующая быстроту изменения вектора скорости \\
\end{tblr}

\subsection{Законы Ньютона}

\subsubsection{Первый закон}
Существуют такие системы отсчета, называемые инерциальными, относительно которых тела сохраняют состояние покоя или равномерного прямолинейного движения, если на них не действуют силы или действие сил скомпенсировано:
\[ \sum \vec{F} = 0 \]

\subsubsection{Второй закон}
В инерциальной системе отсчета ускорение тела прямо пропорционально равнодействующей всех сил, приложенных к телу, и обратно пропорционально его массе:
\[
    \vec{F} = m\vec{a}
\]

\subsubsection{Третий закон}
Силы взаимодействия двух материальных точек равны по величине, противоположно направлены, и действуют вдоль прямой, соединяющей эти материальные точки:
\[
    \vec{F}_{12} = -\vec{F}_{21}
\]

\subsubsection{Выстрел под углом}

Тело бросают с начальной скоростью \( v_0 \) под углом \( \alpha \) к горизонту. Ускорение свободного падения равно \( g \), сопротивлением воздуха пренебречь.
Найти полный путь \( S(\alpha) \), пройденный телом от точки броска до момента падения и угол \( \alpha_{\max} \), при котором этот путь принимает максимальное значение:

\begin{center}
\begin{tikzpicture}
    % Оси
    \draw[->] (-0.5,0) -- (4,0) node[right] {$x$};
    \draw[->] (0,-0.5) -- (0,2.5) node[above] {$y$};
    % Траектория (толстая черная линия для ч/б печати)
    \draw[thick, dashed, domain=0:3.5, smooth, variable=\x] plot ({\x}, {2.2*\x - 0.65*\x*\x});
    % Вектор скорости (толстый черный)
    \draw[thick, -{Stealth}] (0,0) -- (0.8, 1.76) node[above right] {$\vec{v}_0$};
    \draw (0.4,0) arc (0:65:0.4) node[midway, right, xshift=2pt] {$\alpha$};
    % Обозначение пути
    \node at (2.4, 2.1) {Путь $S$};
\end{tikzpicture}
\end{center}

Зависимость координат от времени $t$:
\[ x(t) = v_0 \cos \alpha \cdot t, \quad y(t) = v_0 \sin \alpha \cdot t - \frac{gt^2}{2} \]

Время полета $T$ находится из условия $y(T) = 0$:
\[ T = \frac{2v_0 \sin \alpha}{g} \]

Путь $S$ --- длина дуги траектории, которая вычисляется через интеграл:
\[ S = \int_0^T \sqrt{\left(\frac{dx}{dt}\right)^2 + \left(\frac{dy}{dt}\right)^2} dt \]

Находим компоненты скорости:
\[ v_x = v_0 \cos \alpha, \quad v_y = v_0 \sin \alpha - gt \]

Тогда интеграл принимает вид:
\[ S = \int_0^{T} \sqrt{v_0^2 \cos^2 \alpha + (v_0 \sin \alpha - gt)^2} dt \]

Сделаем замену переменной $u = v_0 \sin \alpha - gt$, тогда $du = -g \, dt$. Пределы интегрирования: от $v_0 \sin \alpha$ до $-v_0 \sin \alpha$.
\[ S = \frac{1}{g} \int_{-v_0 \sin \alpha}^{v_0 \sin \alpha} \sqrt{v_0^2 \cos^2 \alpha + u^2} du \]

Используя табличный интеграл 
\[
\int \sqrt{a^2 + u^2}\, du = \frac{1}{2} \left( u\sqrt{a^2+u^2} + a^2 \ln|u + \sqrt{a^2+u^2}| \right),
\]
получаем выражение для длины дуги траектории:
\[
\boxed{S(\alpha) = \frac{v_0^2}{g} \left[ \sin \alpha + \cos^2 \alpha \cdot \ln \left( \frac{1 + \sin \alpha}{\cos \alpha} \right) \right]}
\]

Для поиска максимума берём производную $S(\alpha)$ по $\alpha$ и приравниваем к нулю. Обозначим $L(\alpha) = \ln\left( \frac{1+\sin\alpha}{\cos\alpha} \right)$. Тогда
\[
S'(\alpha) = \frac{v_0^2}{g} \cdot \frac{d}{d\alpha}\left[ \sin\alpha + \cos^2\alpha \cdot L(\alpha) \right].
\]

Вычислим производную:
\begin{align*}
\frac{d}{d\alpha}\big[\sin\alpha\big] &= \cos\alpha, \\
\frac{d}{d\alpha}\big[\cos^2\alpha \cdot L(\alpha)\big] &= -2\cos\alpha\sin\alpha \cdot L(\alpha) + \cos^2\alpha \cdot L'(\alpha).
\end{align*}

Находим $L'(\alpha)$:
\[
L(\alpha) = \ln(1+\sin\alpha) - \ln(\cos\alpha),
\]
\[
L'(\alpha) = \frac{\cos\alpha}{1+\sin\alpha} + \frac{\sin\alpha}{\cos\alpha} 
= \frac{\cos^2\alpha + \sin\alpha(1+\sin\alpha)}{(1+\sin\alpha)\cos\alpha}
= \frac{1 + \sin\alpha}{(1+\sin\alpha)\cos\alpha} = \frac{1}{\cos\alpha}.
\]

Следовательно,
\[
\cos^2\alpha \cdot L'(\alpha) = \cos\alpha.
\]

Подставляя, получаем:
\[
S'(\alpha) = \frac{v_0^2}{g} \left[ \cos\alpha - 2\cos\alpha\sin\alpha \cdot L(\alpha) + \cos\alpha \right]
= \frac{v_0^2}{g} \cdot 2\cos\alpha \left[ 1 - \sin\alpha \cdot L(\alpha) \right].
\]

Приравнивая производную к нулю при $\alpha \in (0, \pi/2)$ ($\cos\alpha \neq 0$), приходим к уравнению:
\[
1 - \sin\alpha \cdot \ln \left( \frac{1 + \sin \alpha}{\cos \alpha} \right) = 0,
\]
или
\[
\boxed{\sin\alpha \cdot \ln \left( \frac{1 + \sin \alpha}{\cos \alpha} \right) = 1}
\]

Численное решение этого трансцендентного уравнения даёт:
\[
\boxed{\alpha_{\max} \approx 56.45^\circ \quad (\approx 0.985 \text{ рад})}
\]

\subsection{Космические скорости}

\subsubsection{Первая космическая скорость}
Это минимальная скорость, которую необходимо сообщить телу, чтобы оно стало спутником планеты, движущимся по круговой орбите вблизи её поверхности.

На спутник массой $m$, движущийся по круговой орбите радиуса $r = R + h$ вокруг планеты массой $M$, действует только сила гравитационного притяжения. Она сообщает телу центростремительное ускорение.

Согласно II закону Ньютона:
\[
    G \frac{Mm}{r^2} = m \frac{v_1^2}{r}\iff v_1 = \sqrt{\frac{GM}{r}}
\]

Вблизи поверхности Земли ($h \to 0, r \approx R$):
\[
    v_1 = \sqrt{\frac{GM}{R}} = \sqrt{gR} \approx 7.9 \text{ км/с}
\]
где $g = \frac{GM}{R^2}$ — ускорение свободного падения у поверхности.

\subsubsection{Вторая космическая скорость}
Это минимальная скорость, которую нужно сообщить телу, чтобы оно преодолело гравитационное притяжение планеты и ушло в бесконечность (движение по параболе).

Для ухода в бесконечность полная механическая энергия тела должна быть неотрицательной ($E \ge 0$), так как на бесконечности $U \to 0$ и $K \ge 0$.

Запишем закон сохранения энергии для граничного случая ($E_{\infty} = 0$):
\[
    \frac{mv_2^2}{2} - G\frac{Mm}{R} = 0
\]
Отсюда выражаем скорость:
\[
    v_2 = \sqrt{\frac{2GM}{R}} = \sqrt{2} v_1 \approx 11.2 \text{ км/с}
\]

\begin{tblr}{
    colspec = {X[2,c]X[5,l]},
    vlines,
    hlines,
    row{1} = {font=\bfseries, bg=gray!20},
    cells = {m},
    rowsep = 4pt
}
Скорость & Траектория и физический смысл \\
$v < v_1$ & Тело падает на поверхность планеты (баллистическая траектория, эллипс, пересекающий поверхность) \\
$v = v_1$ & Круговая орбита (спутник) \\
$v_1 < v < v_2$ & Эллиптическая орбита \\
$v = v_2$ & Параболическая траектория (тело уходит от планеты, но скорость на бесконечности равна 0) \\
$v > v_2$ & Гиперболическая траектория (тело покидает систему навсегда) \\
\end{tblr}